\NSPage{031}%
\begin{EESegment}{EE-TGT-P031-001}
So the quantities \NSi{NS-F-P031-001}{P_c,J_c} are the scalar products of
\NSi{NS-F-P031-002}{p_s} with the two orthogonal, unnormalized vectors
\NSi{NS-F-P031-003}{(1,t_s)} and \NSi{NS-F-P031-004}{(-t_s,1)}, respectively.
Also, \NSi{NS-F-P031-005}{\mathfrak{s}=a(1,t_s)}. Define
\end{EESegment}
\NSFormula{NS-F-P031-006}{display}{4.21}
\begin{equation}
  \mathcal{U}(P_c,J_c)=P_c+\frac{J_c^2}{4}-|J_c|\sqrt{\frac{P_c-2}{2}+\frac{J_c^2}{16}},
  \qquad P_c>2,\quad v_s<\mathcal{U}(P_c,J_c).
  \tag{4.21}\label{eq:4.21}
\end{equation}
\begin{EESegment}{EE-TGT-P031-002}
When \NSi{NS-F-P031-007}{a>0}, the two inequalities in \eqref{eq:4.21} are
called the \emph{relaxed cone condition}. The \emph{admissible stress cone
condition} asks for those same two inequalities and also requires
\NSi{NS-F-P031-008}{v_s>2}. The relaxed condition is used while the profile
on the axis is being joined to the outer profile. Viscous waves need the
extra inequality. Once the admissible condition holds,
Proposition~\ref{prop:7.5} produces positive squared amplitudes whose wave
covariances give the required stress. The next lemma gives an equivalent way
to check the condition and a sufficient condition used in the outer
construction of Proposition~\ref{prop:A.4}.
\end{EESegment}

\begin{lemma}\label{lem:4.5}
\begin{EESegment}{EE-TGT-P031-003}
Let \NSi{NS-F-P031-009}{a>0}, \NSi{NS-F-P031-010}{b_s\in\mathbb{R}}, and
\NSi{NS-F-P031-011}{p_s=(p_{s,1},p_{s,2})\in\mathbb{R}^2}. Define
\NSi{NS-F-P031-012}{t_s,v_s,P_c,J_c} by \eqref{eq:4.20}. If
\NSi{NS-F-P031-013}{v_s>2}, then the admissible stress cone condition, which
means the inequalities in \eqref{eq:4.21} together with
\NSi{NS-F-P031-014}{v_s>2}, is equivalent to
\end{EESegment}
\NSFormula{NS-F-P031-015}{display}{4.22}
\begin{equation}
  P_c>v_s,\qquad (v_s-2)J_c^2<2(P_c-v_s)^2.
  \tag{4.22}\label{eq:4.22}
\end{equation}

\begin{EESegment}{EE-TGT-P031-004}
Now take any nonempty compact set
\end{EESegment}
\NSFormula{NS-F-P031-016}{display}{none}
\[
  K\subset\{(a,b_s,w)\in\mathbb{R}^3:a>0\}
\]
\begin{EESegment}{EE-TGT-P031-005}
such that every \NSi{NS-F-P031-017}{(a,b_s,w)\in K} satisfies
\end{EESegment}
\NSFormula{NS-F-P031-018}{display}{none}
\[
  a-b_sw>0,\qquad 2b_sw+\frac{b_s^2}{a}+(a-2)w^2<2,
\]
\begin{EESegment}{EE-TGT-P031-006}
There is then a number \NSi{NS-F-P031-019}{P_K>0} with the following
property. For every \NSi{NS-F-P031-020}{(a,b_s,w)\in K} and every
\NSi{NS-F-P031-021}{p_{s,1}\geq P_K}, set
\NSi{NS-F-P031-022}{p_{s,2}=wp_{s,1}}. These data satisfy the relaxed cone
condition \eqref{eq:4.21}. If \NSi{NS-F-P031-023}{v_s>2} as well, they
satisfy the admissible stress cone condition.
\end{EESegment}
\end{lemma}

\begin{proof}
\begin{EESegment}{EE-TGT-P031-007}
Either set of conditions in the first claim gives
\NSi{NS-F-P031-024}{P_c>2}. On this range, the polynomial
\NSi{NS-F-P031-025}{2(P_c-v)^2-(v-2)J_c^2} in
\NSi{NS-F-P031-026}{v} has roots
\end{EESegment}
\NSFormula{NS-F-P031-027}{display}{none}
\[
  v_{\pm}=P_c+\frac{J_c^2}{4}\pm |J_c|\sqrt{\frac{P_c-2}{2}+\frac{J_c^2}{16}}.
\]
\begin{EESegment}{EE-TGT-P031-008}
When \NSi{NS-F-P031-028}{P_c>2}, the lower root satisfies
\NSi{NS-F-P031-029}{2<v_-\leq P_c}. So on
\NSi{NS-F-P031-030}{2<v<P_c}, the polynomial is positive exactly when
\NSi{NS-F-P031-031}{v<v_-}. Since
\NSi{NS-F-P031-032}{\mathcal{U}(P_c,J_c)=v_-}, this proves the first claim.
\end{EESegment}

\begin{EESegment}{EE-TGT-P031-009}
For the second claim, put
\end{EESegment}
\NSFormula{NS-F-P031-033}{display}{none}
\[
  c=1-\frac{b_sw}{a},\qquad j=w+\frac{b_s}{a},\qquad v=a+\frac{b_s^2}{a}=v_s.
\]
\begin{EESegment}{EE-TGT-P031-010}
Then \NSi{NS-F-P031-034}{P_c=p_{s,1}c,\ J_c=p_{s,1}j}, and
\end{EESegment}
\NSFormula{NS-F-P031-035}{display}{none}
\[
  (v_s-2)\left(w+\frac{b_s}{a}\right)^2-2\left(1-\frac{b_sw}{a}\right)^2
  =\left(1+\frac{b_s^2}{a^2}\right)\left((a-2)w^2+2b_sw+\frac{b_s^2}{a}-2\right).
\]
\begin{EESegment}{EE-TGT-P031-011}
The assumptions say that \NSi{NS-F-P031-036}{c>0} and
\NSi{NS-F-P031-037}{G:=2c^2-(v-2)j^2>0} everywhere on
\NSi{NS-F-P031-038}{K}. Because this set is compact, there are constants
\NSi{NS-F-P031-039}{c_K,\gamma_K>0} and
\NSi{NS-F-P031-040}{B_K\geq 1} such that
\end{EESegment}
\NSFormula{NS-F-P031-041}{display}{none}
\[
  c\geq c_K,\qquad G\geq\gamma_K,\qquad v\leq B_K,\qquad cv\leq B_K\quad\hbox{on }K.
\]
\begin{EESegment}{EE-TGT-P031-012}
Choose
\end{EESegment}
\NSFormula{NS-F-P031-042}{display}{none}
\[
  P_K=\max\left\{\frac{B_K+2}{c_K},\frac{8B_K}{\gamma_K}\right\}.
\]
\NSPage{032}%
\begin{EESegment}{EE-TGT-P032-001}
For every \NSi{NS-F-P032-001}{p_{s,1}\geq P_K}, this choice gives
\NSi{NS-F-P032-002}{P_c\geq B_K+2>\max\{2,v\}} and
\end{EESegment}
\NSFormula{NS-F-P032-003}{display}{none}
\[
  \frac{2(P_c-v)^2-(v-2)J_c^2}{p_{s,1}^2}
  =G-\frac{4cv}{p_{s,1}}+\frac{2v^2}{p_{s,1}^2}
  \geq\frac{\gamma_K}{2}>0.
\]
\begin{EESegment}{EE-TGT-P032-002}
If \NSi{NS-F-P032-004}{v>2}, the first claim gives the relaxed condition.
If \NSi{NS-F-P032-005}{v\leq 2}, the relaxed condition follows instead from
\NSi{NS-F-P032-006}{P_c>2} and
\NSi{NS-F-P032-007}{\mathcal{U}(P_c,J_c)>2}. So the same threshold works
everywhere on \NSi{NS-F-P032-008}{K}, including the points where
\NSi{NS-F-P032-009}{v_s=2}.
\end{EESegment}
\end{proof}

\begin{EESegment}{EE-TGT-P032-003}
In stress coordinates, the cone condition is
\end{EESegment}
\NSFormula{NS-F-P032-010}{display}{4.23}
\begin{equation}
  \mathcal{T}_{0,\theta}+t_s\mathcal{T}_{0,z}=F(P_c-v_s),\qquad
  \mathcal{T}_{0,z}-t_s\mathcal{T}_{0,\theta}=FJ_c.
  \tag{4.23}\label{eq:4.23}
\end{equation}
\begin{EESegment}{EE-TGT-P032-004}
Since \NSi{NS-F-P032-011}{F>0}, when
\NSi{NS-F-P032-012}{v_s>2} the admissible stress cone condition says exactly
that
\end{EESegment}
\NSFormula{NS-F-P032-013}{display}{none}
\[
  \mathcal{T}_{0,\theta}+t_s\mathcal{T}_{0,z}>0,
  \qquad
  (v_s-2)(\mathcal{T}_{0,z}-t_s\mathcal{T}_{0,\theta})^2
  <2(\mathcal{T}_{0,\theta}+t_s\mathcal{T}_{0,z})^2.
\]
\begin{EESegment}{EE-TGT-P032-005}
These inequalities are homogeneous in
\NSi{NS-F-P032-014}{\mathcal{T}_0}, so multiplying the stress by a positive
number does not change them. The stress can tend to zero at an edge of the
annulus while its unit direction still satisfies the cone inequalities with
a strict margin. Theorem~\ref{thm:4.6}(iii) states that boundary condition
precisely. Proposition~\ref{prop:7.5} writes every nonzero stress of this
kind as a positive combination of the allowed wave covariances.
Differentiating that representation gives the signed changes in amplitude
used for the higher-order stress corrections in
Proposition~\ref{prop:7.6}.
\end{EESegment}

\begin{EESegment}{EE-TGT-P032-006}
\subsection{Properties of the leading profile}\label{sec:4.4}

The profiles \NSi{NS-F-P032-015}{E,U} must meet both the radial matching
conditions of Lemma~\ref{lem:4.4} and the stress cone inequalities of
Lemma~\ref{lem:4.5}. At each edge of the annulus, the stress
\NSi{NS-F-P032-016}{\mathcal{T}_0} must go smoothly to zero while its direction
stays inside the cone. Theorem~\ref{thm:4.6} puts fixed margins and bounds on these
edge conditions. It also supplies regular behavior at the axis, the heat-flow exterior,
and the moment identities needed in Sections~\ref{sec:2.1} and~\ref{sec:3},
and it keeps aside the intervals later used for the background and mean
corrections in \eqref{eq:4.30}.
\end{EESegment}

\begin{EESegment}{EE-TGT-P032-007}
The edge estimates use the names \emph{inner collar} and \emph{outer collar}
for fixed neighborhoods of the two annular edges in profile coordinates. If
\NSi{NS-F-P032-017}{0<X_a<X_b} and
\NSi{NS-F-P032-018}{0<\varepsilon<\log(X_b/X_a)}, these collars are the
rectangles
\end{EESegment}
\NSFormula{NS-F-P032-019}{display}{4.24}
\begin{equation}
\begin{aligned}
  \mathcal{C}_a(\varepsilon)&=[X_a,X_ae^\varepsilon]\times[-1,1],\\
  \mathcal{C}_b(\varepsilon)&=[X_be^{-\varepsilon},X_b]\times[-1,1].
\end{aligned}
\tag{4.24}\label{eq:4.24}
\end{equation}
\begin{EESegment}{EE-TGT-P032-008}
At a fixed \NSi{NS-F-P032-020}{\eta}, the inner collar has
\NSi{NS-F-P032-021}{0\leq\log(X/X_a)\leq\varepsilon}, while the outer collar
has \NSi{NS-F-P032-022}{0\leq\log(X_b/X)\leq\varepsilon}. Their widths do
not depend on \NSi{NS-F-P032-023}{\eta} or the physical variable
\NSi{NS-F-P032-024}{q}. A smaller collar means the same kind of set with a
smaller positive \NSi{NS-F-P032-025}{\varepsilon}. If a condition is meant
only for points inside the annulus, the collar is intersected with
\NSi{NS-F-P032-026}{\{X_a<X<X_b\}}.
\end{EESegment}

\begin{EESegment}{EE-TGT-P032-009}
A smooth scalar or vector \NSi{NS-F-P032-027}{g} is \emph{flat at the radial
edge} \NSi{NS-F-P032-028}{X_e\in\{X_a,X_b\}} when every one of its radial
derivatives, including the zeroth, vanishes there component by component,
\end{EESegment}
\NSFormula{NS-F-P032-029}{display}{none}
\[
  \partial_X^jg(X_e,\eta)=0\quad\hbox{for every }j\geq0\hbox{ and }\eta\in[-1,1].
\]
\begin{EESegment}{EE-TGT-P032-010}
The stress will be flat at both edges. Wherever the stress is nonzero, its normalized direction is
found by dividing the stress by its length, so it is the unit vector that points in the same direction
as the stress. The theorem below states that the constructed stress is nonzero on the marked open
annulus \NSi{NS-F-P032-030}{X_a<X<X_b}, so its normalized direction is defined there. It also
extends that direction to the closed collars with the positive margin stated in
Theorem~\ref{thm:4.6}(iii). The construction results used in that proof come after the theorem.
\end{EESegment}

\begin{theorem}\label{thm:4.6}
\begin{EESegment}{EE-TGT-P032-011}
There are fixed numbers
\NSi{NS-F-P032-031}{h\in(0,1/100),\ \lambda>0,\ \mathsf C>1} and
\NSi{NS-F-P032-032}{0<X_a<X_b}, together with profiles
\NSi{NS-F-P032-033}{E,U,\Pi} for \NSi{NS-F-P032-034}{X\geq0} and
\NSi{NS-F-P032-035}{-1\leq\eta\leq1}, with the properties listed below. The
quantities
\NSi{NS-F-P032-036}{F,H,V_0,Q_s,N_s,p_s,a,b_s\text{ and }\mathcal{T}_0}
are defined by Equations~\eqref{eq:4.7} to~\eqref{eq:4.11}. Wherever
\NSi{NS-F-P032-037}{a>0}, the quantities \NSi{NS-F-P032-038}{t_s,v_s} are
defined by \eqref{eq:4.20}. The symbols \NSi{NS-F-P032-039}{A,d} come from
\eqref{eq:4.1}. Here the closed annulus means the profile rectangle
\NSi{NS-F-P032-040}{[X_a,X_b]\times[-1,1]} in
\NSi{NS-F-P032-041}{(X,\eta)}. Every claim below holds for every
\NSi{NS-F-P032-042}{\eta\in[-1,1]}.
\end{EESegment}
\begin{enumerate}[label=(\roman*)]
\NSPage{033}%
\item
\begin{EESegment}{EE-TGT-P033-001}
For every finite \NSi{NS-F-P033-001}{R>0}, the scalars
\NSi{NS-F-P033-002}{F=\phi/\mathsf C,U,\Pi,V_0/X} are smooth on
\NSi{NS-F-P033-003}{[0,R]\times[-1,1]}. Every fixed mixed
\NSi{NS-F-P033-004}{X,\eta} derivative has a uniform bound there, with
one-sided derivatives at the boundary. Also, \NSi{NS-F-P033-005}{\phi>0},
and \NSi{NS-F-P033-006}{E=\sqrt{2X}F>0} whenever
\NSi{NS-F-P033-007}{X>0}. By Equations~\eqref{eq:4.4}--\eqref{eq:4.5}, these
factors give a smooth Cartesian field across the axis.
\end{EESegment}

\begin{EESegment}{EE-TGT-P033-002}
There is one fixed \NSi{NS-F-P033-008}{X_{\mathrm{an}}\in(X_a,X_b)} with the
following property. On
\NSi{NS-F-P033-009}{[0,X_{\mathrm{an}}]\times[-1,1]}, these scalars and each
fixed radial derivative are analytic in \NSi{NS-F-P033-010}{\eta} on one
common complex neighborhood of \NSi{NS-F-P033-011}{[-1,1]}. The inner collar
\NSi{NS-F-P033-012}{[X_a,X_{\mathrm{an}}]\times[-1,1]}, understood as in
\eqref{eq:4.24}, has the positive directional margin stated in part~(iii).
The pressure is normalized to vanish at radial infinity, so
\end{EESegment}
\NSFormula{NS-F-P033-013}{display}{4.25}
\begin{equation}
  \Pi(X,\eta)=-\int_X^\infty\frac{E(x,\eta)^2}{2x}\,dx.
  \tag{4.25}\label{eq:4.25}
\end{equation}

\item
\begin{EESegment}{EE-TGT-P033-003}
The radial pressure balance
\NSi{NS-F-P033-014}{\Pi_X=E^2/(2X)} in \eqref{eq:4.7}, with its smooth extension
\NSi{NS-F-P033-015}{\Pi_X=F^2} at \NSi{NS-F-P033-016}{X=0}, holds exactly.
So do the tangential residual identities in Proposition~\ref{prop:4.2}, with
physical stress \NSi{NS-F-P033-017}{q^{-A-1/2}\mathcal{T}_0}. The stress
profile \NSi{NS-F-P033-018}{\mathcal{T}_0} is zero when
\NSi{NS-F-P033-019}{0\leq X\leq X_a\text{ and }X\geq X_b}, and it is nonzero
at every point with \NSi{NS-F-P033-020}{X_a<X<X_b}. On
\NSi{NS-F-P033-021}{0\leq X\leq X_a}, the profiles solve \eqref{eq:4.13}.
\end{EESegment}

\item
\begin{EESegment}{EE-TGT-P033-004}
Each of \NSi{NS-F-P033-022}{F,a,v_s-2} has a positive lower bound on the
closed annulus \NSi{NS-F-P033-023}{[X_a,X_b]\times[-1,1]}. The unit
direction \NSi{NS-F-P033-024}{n=\mathcal{T}_0/|\mathcal{T}_0|} starts out
defined on the radial interior
\NSi{NS-F-P033-025}{(X_a,X_b)\times[-1,1]}, and it extends smoothly to the
whole closed annulus, including the boundaries
\NSi{NS-F-P033-026}{X=X_a,X_b}. There is one fixed
\NSi{NS-F-P033-027}{\kappa\in(0,2)} such that everywhere on this rectangle,
\end{EESegment}
\NSFormula{NS-F-P033-028}{display}{4.26}
\begin{equation}
  n_\theta+t_sn_z\geq\kappa,\qquad
  (v_s-2)(n_z-t_sn_\theta)^2\leq(2-\kappa)(n_\theta+t_sn_z)^2.
  \tag{4.26}\label{eq:4.26}
\end{equation}
\begin{EESegment}{EE-TGT-P033-005}
For each \NSi{NS-F-P033-029}{\eta}, the vector
\NSi{NS-F-P033-030}{n(X_a,\eta)\text{ is parallel to }(a(X_a,\eta),-b_s(X_a,\eta))}.
At \NSi{NS-F-P033-031}{X=X_b},
\NSi{NS-F-P033-032}{n(X_b,\eta)=(1,0)} and
\NSi{NS-F-P033-033}{b_s(X_b,\eta)=0}.
\end{EESegment}

\item
\begin{EESegment}{EE-TGT-P033-006}
There is a smooth weight \NSi{NS-F-P033-034}{\zeta(X)}. It is positive when
\NSi{NS-F-P033-035}{X_a<X<X_b\text{ and zero for }X\leq X_a\text{ and }X\geq X_b}.
Every radial derivative of \NSi{NS-F-P033-036}{\zeta} vanishes at
\NSi{NS-F-P033-037}{X_a,X_b}. In the radial interior, put
\NSi{NS-F-P033-038}{\delta=\min\{1,\log(X/X_a),\log(X_b/X)\}}. There is a
constant \NSi{NS-F-P033-039}{c>0} that does not depend on
\NSi{NS-F-P033-040}{\alpha}. For each fixed multi-index
\NSi{NS-F-P033-041}{\alpha=(\alpha_1,\alpha_2)}, there are finite constants
\NSi{NS-F-P033-042}{C_\alpha,m_\alpha} such that, with
\NSi{NS-F-P033-043}{\partial^\alpha=\partial_X^{\alpha_1}\partial_\eta^{\alpha_2}},
the following bounds hold throughout
\NSi{NS-F-P033-044}{(X_a,X_b)\times[-1,1]}.
\end{EESegment}
\NSFormula{NS-F-P033-045}{display}{4.27}
\begin{equation}
  |\mathcal{T}_0|\geq c\zeta,\qquad
  |\partial^\alpha\mathcal{T}_0|\leq C_\alpha\zeta\delta^{-m_\alpha}.
  \tag{4.27}\label{eq:4.27}
\end{equation}

\item
\begin{EESegment}{EE-TGT-P033-007}
The limits \NSi{NS-F-P033-046}{M(\infty,\eta),J(\infty,\eta),S(\infty,\eta)}
of the cumulative integrals in \eqref{eq:4.15} exist. The angular moment is
made convergent by subtracting the exterior power
\NSi{NS-F-P033-047}{H_{\mathrm{pow}}} below from
\NSi{NS-F-P033-048}{H} before integrating. These four convergent quantities
satisfy
\end{EESegment}
\NSFormula{NS-F-P033-049}{display}{4.28}
\begin{equation}
  M(\infty,\eta)=J(\infty,\eta)=S(\infty,\eta)=0,\qquad
  \int_0^\infty(H-H_{\mathrm{pow}})\,dX=0,\qquad
  H_{\mathrm{pow}}=\sqrt{2X}\,c_\infty X^{-A},
  \tag{4.28}\label{eq:4.28}
\end{equation}
\begin{EESegment}{EE-TGT-P033-008}
where \NSi{NS-F-P033-050}{c_\infty>0} does not depend on
\NSi{NS-F-P033-051}{\eta}. For \NSi{NS-F-P033-052}{X\geq X_b},
\end{EESegment}
\NSFormula{NS-F-P033-053}{display}{4.29}
\begin{equation}
\begin{aligned}
  U=V_0=0,\qquad&E=c_\infty X^{-A}\mathcal{H}(2d/X),\\
  \mathcal{H}(Z)&=\frac{1}{\Gamma(1+h)}\int_0^\infty e^{-v}v^h(1+Zv)^{-h}\,dv.
\end{aligned}
\tag{4.29}\label{eq:4.29}
\end{equation}
\begin{EESegment}{EE-TGT-P033-009}
In this region the physical swirl is the exact radial heat solution
\NSi{NS-F-P033-054}{c_\infty s^{-A}\mathcal{H}(2\tau/s)}. The profile
\NSi{NS-F-P033-055}{\mathcal{H}} is positive and smooth for
\NSi{NS-F-P033-056}{Z\geq0}, including the one-sided endpoint at zero. There
is also a fixed \NSi{NS-F-P033-057}{X_v\in(X_a,X_b)} such that
\NSi{NS-F-P033-058}{U=V_0=0} on all of
\NSi{NS-F-P033-059}{[X_v,\infty)}, including the outer collar
\NSi{NS-F-P033-060}{[X_v,X_b]}.
\end{EESegment}

\item
\begin{EESegment}{EE-TGT-P033-010}
Two fixed, disjoint intervals
\NSi{NS-F-P033-061}{\mathcal{I}_{\mathrm{pos}}\text{ and }\mathcal{I}_{\mathrm{mean}}}
are left available for later use. Lemma~\ref{lem:5.2} uses the first for
positive-order background corrections, and Lemma~\ref{lem:8.7} uses the
second for corrections to the averaged flow,
\end{EESegment}
\NSPage{034}%
\begin{EESegment}{EE-TGT-P034-001}
with \NSi{NS-F-P034-001}{\sup\mathcal{I}_{\mathrm{pos}}<\inf\mathcal{I}_{\mathrm{mean}}}.
Here positive order means a relative power
\NSi{NS-F-P034-002}{q^{2nh},\ n\geq1} in the background expansion
\eqref{eq:5.1}. Both intervals lie in
\NSi{NS-F-P034-003}{(X_a,X_v)}. For every \NSi{NS-F-P034-004}{X} in either
interval,
\end{EESegment}
\NSFormula{NS-F-P034-005}{display}{4.30}
\begin{equation}
  U=0,\qquad E=c_{\mathrm{patch}}(1+\eta^2)^{-1}X^{-1/2-\lambda},
  \tag{4.30}\label{eq:4.30}
\end{equation}
\begin{EESegment}{EE-TGT-P034-002}
where the positive constant \NSi{NS-F-P034-006}{c_{\mathrm{patch}}} does not
depend on \NSi{NS-F-P034-007}{X,\eta}. Building the leading profile does not
change either interval.
\end{EESegment}
\end{enumerate}

\begin{EESegment}{EE-TGT-P034-003}
Every constant in this theorem depends only on the fixed choices used to
build the profile and, for a derivative bound, on the order of that
derivative. None depends on the physical variable
\NSi{NS-F-P034-008}{q}, on any later dyadic band, or on any later correction
stage.
\end{EESegment}
\end{theorem}

\begin{EESegment}{EE-TGT-P034-004}
\subsection{Construction results used in the proof}\label{sec:4.5}

Building the profile takes three choices that have to work together. The
outer flow fixes the pressure on the axis, the regular inner flow matches the outer flow's
five cumulative integrals, and the shear on the interval between them
satisfies the admissible cone condition. Lemma~\ref{lem:4.7} first solves the
finite-dimensional moment problem. Lemma~\ref{lem:4.8} and
Proposition~\ref{prop:4.10} then build the outer and inner flows, and
Lemma~\ref{lem:4.11} builds the periodic shear. These results provide the
pieces used in the proof of Theorem~\ref{thm:4.6}. The appendices prove each
construction in full at the references given below.
\end{EESegment}

\begin{EESegment}{EE-TGT-P034-005}
In these statements, a \emph{parameter derivative} means a derivative with
respect to \NSi{NS-F-P034-009}{\eta}, with every construction constant held
fixed. For a function on \NSi{NS-F-P034-010}{[-1,1]}, write
\NSi{NS-F-P034-011}{\|f\|_{C^k}=\max_{0\leq j\leq k}\sup_\eta|\partial_\eta^jf|}.
For vectors, use the largest component norm. Only finitely many of these
norms have to be small when the parameters are chosen. After the parameters
have been fixed, every other fixed derivative order has a finite bound.
\end{EESegment}

\begin{lemma}\label{lem:4.7}
\begin{EESegment}{EE-TGT-P034-006}
Let \NSi{NS-F-P034-012}{\alpha_1,\ldots,\alpha_m} be distinct real numbers.
Let \NSi{NS-F-P034-013}{\beta_1,\ldots,\beta_m} be nonnegative, nonzero
smooth functions supported in ordered, disjoint compact intervals of
\NSi{NS-F-P034-014}{(0,\infty)}. Then the matrix
\end{EESegment}
\NSFormula{NS-F-P034-015}{display}{none}
\[
  B_{ij}=\int_0^\infty x^{\alpha_i}\beta_j(x)\,dx
\]
\begin{EESegment}{EE-TGT-P034-007}
is invertible. Its inverse and every fixed parameter derivative stay bounded
for any smooth compact family that keeps these assumptions true.
\end{EESegment}

\begin{EESegment}{EE-TGT-P034-008}
There is also a more general statement. Let
\NSi{NS-F-P034-016}{B(\eta)} be a smooth invertible matrix on
\NSi{NS-F-P034-017}{[-1,1]}, let
\NSi{NS-F-P034-018}{\mathcal Q_\eta} be a smooth bilinear map, and let
\NSi{NS-F-P034-019}{d\in C^\infty([-1,1];\mathbb{R}^m)}. Fix an integer
\NSi{NS-F-P034-020}{k\geq0}. Let \NSi{NS-F-P034-021}{\nu_j} be the norm of multiplication by
\NSi{NS-F-P034-022}{B^{-1}} on \NSi{NS-F-P034-023}{C^j}, and let
\NSi{NS-F-P034-024}{q_j} be the bilinear norm of
\NSi{NS-F-P034-025}{\mathcal Q} on \NSi{NS-F-P034-026}{C^j}, for
\NSi{NS-F-P034-027}{j=0,k}. If
\end{EESegment}
\NSFormula{NS-F-P034-028}{display}{none}
\[
  8\nu_j^2q_j\|d\|_{C^j}\leq1\qquad(j=0,k),
\]
\begin{EESegment}{EE-TGT-P034-009}
then the equation
\end{EESegment}
\NSFormula{NS-F-P034-029}{display}{none}
\[
  B(\eta)c(\eta)+\mathcal Q_\eta(c(\eta),c(\eta))=d(\eta)
\]
\begin{EESegment}{EE-TGT-P034-010}
has a smooth solution. It is the only solution with
\NSi{NS-F-P034-030}{\|c\|_{C^0}\leq2\nu_0\|d\|_{C^0}}, and it satisfies
\NSi{NS-F-P034-031}{\|c\|_{C^k}\leq2\nu_k\|d\|_{C^k}}. When
\NSi{NS-F-P034-032}{\mathcal Q=0}, no smallness condition is needed and the
solution is \NSi{NS-F-P034-033}{c=B^{-1}d}.
\end{EESegment}
\end{lemma}

\begin{proof}
\begin{EESegment}{EE-TGT-P034-011}
These two claims are Lemmas~\ref{lem:A.1} and~\ref{lem:A.2}. For the first,
multilinearity gives
\end{EESegment}
\NSFormula{NS-F-P034-034}{display}{none}
\[
  \det B=\int\det[x_j^{\alpha_i}]\prod_j\beta_j(x_j)\,dx_1\cdots dx_m.
\]
\begin{EESegment}{EE-TGT-P034-012}
When \NSi{NS-F-P034-035}{x_1<\cdots<x_m}, the determinant inside this
integral has one fixed nonzero sign. To see why, any nonzero linear
combination of \NSi{NS-F-P034-036}{m} distinct powers has at most
\NSi{NS-F-P034-037}{m-1} positive zeros, by induction and Rolle's theorem.
So the integral is nonzero.
\end{EESegment}

\begin{EESegment}{EE-TGT-P034-013}
For the second claim, on the ball of radius
\NSi{NS-F-P034-038}{r_j=2\nu_j\|d\|_{C^j}}, the iteration
\NSi{NS-F-P034-039}{c\mapsto B^{-1}(d-\mathcal Q(c,c))} satisfies
\end{EESegment}
\NSFormula{NS-F-P034-040}{display}{none}
\[
  \|B^{-1}(d-\mathcal Q(c,c))\|_{C^j}\leq r_j/2+\nu_jq_jr_j^2\leq r_j,
  \qquad 2\nu_jq_jr_j\leq\tfrac12.
\]
\NSPage{035}%
\begin{EESegment}{EE-TGT-P035-001}
It is a contraction in \NSi{NS-F-P035-001}{C^0} and
\NSi{NS-F-P035-002}{C^k}. Starting the iteration from zero picks the same
solution in both spaces. The pointwise implicit function theorem then shows
that this solution is smooth.
\end{EESegment}
\end{proof}

\begin{EESegment}{EE-TGT-P035-002}
The exterior is built before the part near the axis. Requiring its pressure
to vanish at infinity fixes the pressure value that the axis problem must
use. The family below contains the heat-flow exterior and keeps several
intervals available, so later changes can restore the required moments.
\end{EESegment}

\begin{lemma}\label{lem:4.8}
\begin{EESegment}{EE-TGT-P035-003}
There are finite choices of the parameters
\NSi{NS-F-P035-003}{M_d,P_\ast,\lambda,h} and a constant
\NSi{NS-F-P035-004}{R_\ast>0} with the following properties. First choose
\NSi{NS-F-P035-005}{M_d}. Then \NSi{NS-F-P035-006}{P_\ast} may depend on
\NSi{NS-F-P035-007}{M_d}; \NSi{NS-F-P035-008}{\lambda} may depend on
\NSi{NS-F-P035-009}{M_d,P_\ast}; and \NSi{NS-F-P035-010}{h} may depend on
\NSi{NS-F-P035-011}{M_d,P_\ast,\lambda}. The choices satisfy
\end{EESegment}
\NSFormula{NS-F-P035-012}{display}{none}
\[
  T_d=e^{M_d}+10,\qquad P_\ast>e^{T_d},\qquad
  0<h<\min\{1/100,\lambda,e^{-T_d}\}.
\]
\begin{EESegment}{EE-TGT-P035-004}
For every \NSi{NS-F-P035-013}{X_R\geq R_\ast}, there are smooth profiles
\NSi{NS-F-P035-014}{(U_o,E_o)} on
\NSi{NS-F-P035-015}{(0,\infty)\times[-1,1]}, with
\NSi{NS-F-P035-016}{E_o>0}, that satisfy the conditions below. Put
\NSi{NS-F-P035-017}{x=X/X_R} and
\NSi{NS-F-P035-018}{f(\eta)=(1+\eta^2)^{-1}}.
\end{EESegment}
\begin{enumerate}[label=(\roman*)]
\item
\begin{EESegment}{EE-TGT-P035-005}
When \NSi{NS-F-P035-019}{0<x\leq1},
\end{EESegment}
\NSFormula{NS-F-P035-020}{display}{none}
\[
  U_o=4\eta,\qquad E_o=P_\ast f(\eta)x^{1/10}.
\]
\begin{EESegment}{EE-TGT-P035-006}
The pressure value
\end{EESegment}
\NSFormula{NS-F-P035-021}{display}{4.31}
\begin{equation}
  \Pi_0(\eta)=-\int_0^\infty\frac{E_o(X,\eta)^2}{2X}\,dX,
  \qquad
  \Pi_o(X,\eta)=\Pi_0(\eta)+C_{p,o}(X,\eta)
  \tag{4.31}\label{eq:4.31}
\end{equation}
\begin{EESegment}{EE-TGT-P035-007}
does not depend on \NSi{NS-F-P035-022}{X_R}. It is analytic on one complex
neighborhood of \NSi{NS-F-P035-023}{[-1,1]}, it is even, and it satisfies
\end{EESegment}
\NSFormula{NS-F-P035-024}{display}{none}
\[
  \Pi_0\leq-\frac52P_\ast^2f^2,\qquad \eta\Pi_0'(\eta)>0\quad(\eta\neq0).
\]
\begin{EESegment}{EE-TGT-P035-008}
For this temporary profile, define the five cumulative integrals by
\eqref{eq:4.15}, \NSi{NS-F-P035-025}{V_{0,o}} by \eqref{eq:4.7}, and
\NSi{NS-F-P035-026}{Q_{s,o},N_{s,o}} by the explicit formulas
\eqref{eq:4.16}. Equations~\eqref{eq:4.11} and~\eqref{eq:4.20} then define
its shear, stress, and cone coordinates. These are the same definitions used
in Lemma~\ref{lem:4.4}. They make sense on
\NSi{NS-F-P035-027}{X>0}, even though this temporary inner branch does not
have to be regular at the physical axis.
\end{EESegment}

\item
\begin{EESegment}{EE-TGT-P035-009}
There are four nonempty, ordered, disjoint intervals
\NSi{NS-F-P035-028}{\mathcal{I}_j=X_R(\alpha_j,\beta_j)} and radii satisfying
\end{EESegment}
\NSFormula{NS-F-P035-029}{display}{none}
\[
  X_R<X_{\mathrm{good}}<\inf\mathcal{I}_1,\qquad
  \sup\mathcal{I}_4<X_v<X_{\mathrm{tail}}<X_b=e^3X_{\mathrm{tail}},
\]
\begin{EESegment}{EE-TGT-P035-010}
and the endpoints of those intervals, after division by
\NSi{NS-F-P035-030}{X_R}, do not depend on
\NSi{NS-F-P035-031}{X_R,\eta}. On all four intervals,
\NSi{NS-F-P035-032}{U_o=0}, and
\end{EESegment}
\NSFormula{NS-F-P035-033}{display}{none}
\[
  E_o=c_{\mathrm{patch}}f(\eta)X^{-1/2-\lambda}\quad\hbox{on }\mathcal{I}_1,\mathcal{I}_3,\mathcal{I}_4,
  \qquad c_{\mathrm{patch}}>0.
\]
\begin{EESegment}{EE-TGT-P035-011}
Only \NSi{NS-F-P035-034}{\mathcal{I}_2} is used for the heat compensation.
The other three intervals remain available for corrections. Also,
\end{EESegment}
\NSFormula{NS-F-P035-035}{display}{none}
\[
  U_o=M_o=J_o=0\qquad(X\geq X_v),
\]
\NSFormula{NS-F-P035-036}{display}{none}
\[
  M_o(\infty)=J_o(\infty)=S_o(\infty)=0,
  \qquad \int_0^\infty(H_o-H_{\mathrm{pow}})\,dX=0,
\]
\begin{EESegment}{EE-TGT-P035-012}
where \NSi{NS-F-P035-037}{H_o=\sqrt{2X}E_o} and
\NSi{NS-F-P035-038}{H_{\mathrm{pow}}=\sqrt{2X}c_\infty X^{-A}}. Both
\NSi{NS-F-P035-039}{c_{\mathrm{patch}},c_\infty>0} are independent of
\NSi{NS-F-P035-040}{X,\eta}.
\end{EESegment}

\item
\begin{EESegment}{EE-TGT-P035-013}
The relaxed cone condition holds on
\NSi{NS-F-P035-041}{[e^{-5}X_R,e^{1/2}X_{\mathrm{tail}}]\times\allowbreak[-1,1]}.
The admissible condition holds on
\NSi{NS-F-P035-042}{[X_{\mathrm{good}},e^{1/2}X_{\mathrm{tail}}]\times\allowbreak[-1,1]}.
On each closed rectangle, every gap that defines the relevant condition has
a positive minimum. These minima may depend on
\NSi{NS-F-P035-043}{X_R}. For \NSi{NS-F-P035-044}{X\geq X_b},
\end{EESegment}
\NSFormula{NS-F-P035-045}{display}{none}
\[
  U_o=V_{0,o}=0,\qquad E_o=c_\infty X^{-A}\mathcal{H}(2d/X),
\]
\begin{EESegment}{EE-TGT-P035-014}
where \NSi{NS-F-P035-046}{\mathcal{H}} is defined in \eqref{eq:4.29}. The
physical swirl
\NSi{NS-F-P035-047}{K=c_\infty(r^2/2)^{-A}\mathcal{H}(4\tau/r^2)} satisfies
\NSi{NS-F-P035-048}{\partial_tK=(\partial_{rr}+r^{-1}\partial_r-r^{-2})K}
and \NSi{NS-F-P035-049}{K_r<0\text{ for }r>0}.
\end{EESegment}
\end{enumerate}
\end{lemma}

\NSPage{036}%
\begin{proof}
\begin{EESegment}{EE-TGT-P036-001}
Proposition~\ref{prop:A.4} and Lemma~\ref{lem:A.5} provide the reference
profiles and the common pressure value in part~(i). The four intervals
defined after \eqref{eq:A.9} are the intervals in part~(ii). The parameter
\NSi{NS-F-P036-001}{M_d} controls the decrease in the axial direction,
\NSi{NS-F-P036-002}{T_d} controls how long that decrease lasts in logarithmic
radius, \NSi{NS-F-P036-003}{P_\ast} is the initial azimuthal amplitude, and
\NSi{NS-F-P036-004}{\lambda,h} are the intermediate and outer exponents.
\end{EESegment}

\begin{EESegment}{EE-TGT-P036-002}
Lemma~\ref{lem:A.6} and Proposition~\ref{prop:A.7} replace the tail by the
heat profile from part~(iii). The correction on
\NSi{NS-F-P036-005}{\mathcal{I}_2} restores the pressure integral from
part~(i), along with the total \NSi{NS-F-P036-006}{S}-moment and the
renormalized angular moment from part~(ii). Once the other parameters have
been fixed, those results keep the cone inequalities from part~(iii) true
for every sufficiently large \NSi{NS-F-P036-007}{X_R}.
\end{EESegment}
\end{proof}

\begin{EESegment}{EE-TGT-P036-003}
The exterior is now prepared with the right moments and heat flow. Its stress
was defined by integrating outward from the axis. To get the conclusions at
the outer edge as well, it still needs a regular inner extension with the
same moments. Once that condition is met, the stress can instead be
integrated backward from infinity, and its limiting direction can be found
explicitly.
\end{EESegment}

\begin{lemma}\label{lem:4.9}
\begin{EESegment}{EE-TGT-P036-004}
For every \NSi{NS-F-P036-008}{X_R\geq R_\ast}, the family in
Lemma~\ref{lem:4.8} can be chosen to have the following extra property.
Suppose a leading field is regular at the axis, agrees with that family when
\NSi{NS-F-P036-009}{X\geq X_{\mathrm{tail}}}, satisfies the four moment
identities in \eqref{eq:4.28} as functions of
\NSi{NS-F-P036-010}{\eta}, and has the canonical pressure \eqref{eq:4.25}.
Then its physical stress \NSi{NS-F-P036-011}{T=q^{-A-1/2}\mathcal T_0}
satisfies
\end{EESegment}
\NSFormula{NS-F-P036-012}{display}{none}
\[
  T_\theta(r)=r^{-2}\int_r^\infty r'^2R_\theta^{(0)}(r')\,dr',
  \qquad
  T_z(r)=r^{-1}\int_r^\infty r'R_z^{(0)}(r')\,dr',
\]
\begin{EESegment}{EE-TGT-P036-005}
so \NSi{NS-F-P036-013}{\mathcal{T}_0=0} whenever
\NSi{NS-F-P036-014}{X\geq X_b}. On
\NSi{NS-F-P036-015}{e^{1/2}X_{\mathrm{tail}}\leq X<X_b}, the admissible cone
condition holds and
\end{EESegment}
\NSFormula{NS-F-P036-016}{display}{none}
\[
  b_s=0,\qquad 2+h<a\leq2+2h,\qquad T_\theta>0,
  \qquad 2-(a-2)(T_z/T_\theta)^2\geq\kappa_o>0.
\]
\begin{EESegment}{EE-TGT-P036-006}
Put \NSi{NS-F-P036-017}{y_b=\log(X_b/X)}. There are smooth functions
\NSi{NS-F-P036-018}{b_\theta,b_z} on a closed outer collar, with
\NSi{NS-F-P036-019}{\inf b_\theta>0}, such that
\end{EESegment}
\NSFormula{NS-F-P036-020}{display}{4.32}
\begin{equation}
  \mathcal{T}_{0,\theta}=e^{-4/y_b^2}y_b^{-3}b_\theta(y_b,\eta),
  \qquad
  \mathcal{T}_{0,z}=e^{-4/y_b^2}y_b^3b_z(y_b,\eta).
  \tag{4.32}\label{eq:4.32}
\end{equation}
\begin{EESegment}{EE-TGT-P036-007}
For this fixed profile, all these claims hold uniformly for
\NSi{NS-F-P036-021}{\eta\in[-1,1]}. Their constants may depend on
\NSi{NS-F-P036-022}{X_R}, but they do not depend on the physical variable
\NSi{NS-F-P036-023}{q}.
\end{EESegment}
\end{lemma}

\begin{proof}
\begin{EESegment}{EE-TGT-P036-008}
These claims are Lemma~\ref{lem:A.8} and
Proposition~\ref{prop:A.10}. The four moment identities make the total
residual integrals cancel. This gives the backward primitives shown above,
which are the integrals from the current radius to infinity. The prepared
terminal profile then gives the cone inequalities and
\eqref{eq:4.32}.
\end{EESegment}
\end{proof}

\begin{EESegment}{EE-TGT-P036-009}
The missing inner extension can now be built. It must be smooth at the axis
and must match the five cumulative integrals of the prepared exterior at a
chosen joining radius. It uses the function
\NSi{NS-F-P036-024}{\Pi_0(\eta)} from Lemma~\ref{lem:4.8} as its pressure
value at \NSi{NS-F-P036-025}{X=0}.
\end{EESegment}

\begin{proposition}\label{prop:4.10}
\begin{EESegment}{EE-TGT-P036-010}
Fix the outer-profile data from Lemma~\ref{lem:4.8} and any prescribed finite
set of \NSi{NS-F-P036-026}{\eta}-derivative orders. It is possible to choose
finite axis parameters
\NSi{NS-F-P036-027}{j_0,\delta_\ast,\sigma_\ast,\Lambda}, a logarithmic
transition length \NSi{NS-F-P036-028}{T_{\mathrm{sh}}}, and an amplitude
threshold \NSi{NS-F-P036-029}{\mathsf C_0} so that every
\NSi{NS-F-P036-030}{\mathsf C\geq \mathsf C_0} allows the following
construction. Set
\end{EESegment}
\NSFormula{NS-F-P036-031}{display}{none}
\[
  X_a=4/\Lambda,\qquad X_R=110(\mathsf C P_\ast)^{10},\qquad
  x=X/X_R,\qquad X_h=X_Re^{-5},\qquad f=(1+\eta^2)^{-1}.
\]
\begin{EESegment}{EE-TGT-P036-011}
After choosing \NSi{NS-F-P036-032}{\mathsf C}, choose a positive activation
width \NSi{NS-F-P036-033}{t_1} in the logarithmic coordinate
\NSi{NS-F-P036-034}{\log(X/X_a)}, a positive shear factor
\NSi{NS-F-P036-035}{\kappa_0}, and the remaining transition widths. There
are profiles \NSi{NS-F-P036-036}{E=\sqrt{2X}\phi/\mathsf C,U} and
\NSi{NS-F-P036-037}{\Pi=\Pi_0+C_p} on
\NSi{NS-F-P036-038}{[0,X_h]\times[-1,1]} with the following properties.
\end{EESegment}
\begin{enumerate}[label=(\roman*)]
\item
\begin{EESegment}{EE-TGT-P036-012}
The scalar \NSi{NS-F-P036-039}{\phi} is positive, and
\NSi{NS-F-P036-040}{\phi,U,\Pi,V_0/X} are smooth on
\NSi{NS-F-P036-041}{[0,X_h]\times[-1,1]}, including at
\NSi{NS-F-P036-042}{X=0}. So Equations~\eqref{eq:4.4}--\eqref{eq:4.5}
give a Cartesian velocity and pressure that are smooth across the
axis. For some
\NSi{NS-F-P036-043}{X_{\mathrm{an}}\in(X_a,X_h)}, these scalars and every
fixed radial derivative are analytic in \NSi{NS-F-P036-044}{\eta} on one
common complex neighborhood throughout
\NSi{NS-F-P036-045}{[0,X_{\mathrm{an}}]\times[-1,1]}. The profiles solve
\eqref{eq:4.13} on \NSi{NS-F-P036-046}{0\leq X\leq X_a}, and
\NSi{NS-F-P036-047}{\mathcal{T}_0=0} there.
\end{EESegment}
